% !TeX root = ../papers/nc-p3.tex
\section{順中心と反対の逆中心による分岐階数十一の下界}
\label{sec:branch-rank-eleven-floor}
\begin{theorem}\label{thm:branch-rank-eleven-floor}
曲面射影のない標準正則窓に純関係がなく片側分岐があれば、反対階数は
$\rho\ge11$ である。$\rho=11$ なら前節の相対二位置で、
順中心 $p$ とその反対位置の逆中心 $d$ を取る。$p$ の後続直線と $d$ の前点直線は
ただ一点 $e$ で交わる。反対切断の両側の行空間について
\[
 \dim W_p=\dim W_d=\dim W_e=4,\quad
 W_p\cap W_d=k w(p,d),\quad W_e=(W_p+W_d)^\perp .
\]
共通元は階数1である。とくに反対階数9と10では純関係のない全ての曲線射影が
一点ファイバーのグラフ判定へ戻る。
\end{theorem}
\begin{proof}
まずどの純関係自由な順中心 $p$ でも後続直線 $L$ の全文字で中央支持は4である。
系\ref{cor:pure-free-branch-rank-nine}で一般点について成立し、
定理\ref{thm:finite-branch-corridors}の $w(x)=\sum K_j(x)\otimes A_j$ の表示で
三つの $K_j$ が独立な点は同じ中央支持を持つ。従属点は順中心であり、その縮約自身が
$L_3$ 型なので中央支持4である。逆向きにも同じである。

二直線列なら命題\ref{prop:two-line-branch-block}から $\rho\ge12$ である。
その他では前定理により順中心 $p$ に反対位置の逆中心 $d$ がある。
$p$ を位置0、$d$ を位置2とし、それぞれの位置1での支持を $S,T$ とする。
どちらも二次元である。$S=T$ なら $p$ の後続直線と $d$ の前点直線が一致し、
一般点はその二定点を結ぶので一つしか動く位置のない列になり矛盾する。
一方 $w(p,d)$ は $w(p)$ の四次元中央支持によって非零であり、
位置1では $S\cap T$ に支持される。よって $S\cap T$ は直線で、
二つの点ファイバーはただ一点 $e$ で交わる。

$W_p=w(p,-)$ と $W_d=w(-,d)$ はどちらも四次元で、それぞれ位置1,3の二次行列の空間である。
共通元は位置1で $S\cap T$ に支持される。$W_p$ は二行三列のHankel四平面であり、
固定した像直線に入る行列は一次元である。例えば非零な左方向で縮約すれば、
二項三次の四係数への隣り合う係数の三本の条件は独立である。
従って共通空間は非零な $w(p,d)$ の直線に一致する。
$W_e$ は $L$ 上なので次元4で、隣接する両中心の行空間を消す。
従って $\rho\ge(4+4-1)+4=11$、等号で直交補の主張も従う。
\end{proof}

\begin{proposition}\label{prop:branch-first-pencil-splitting}
純関係自由な順中心の後続直線では、$W_p^\perp$ の定数束に入る中支持束は
$4\mathcal O(-1)$ であり、余核束の次数は4、階数は $\rho-8$ である。
第二の動く像が有理正規三次なら、余核の各直和因子 $\mathcal O(a)$ は $0\le a\le3$ を満たす。
特に $\rho=11$ では分裂次数は順序を除いて $(1,1,2)$ である。
\end{proposition}
\begin{proof}
位置3の全文字による $w(x,-)$ は直線の各点で単射なので主張の部分束となる。
余核は大域生成され行列式の次数は4である。
第二曲線が三次なら、そこからの偶位置の行像を $W_p$ を法として読むと、
隣接消失により余核束の双対の3捻りに入る切断となる。
これらの切断の全値は定数空間 $W/W_p$ を張る。実際三次曲線の評価点が位置2の全空間を張り、
位置0では全文字を自由に取るためである。
もし $a>3$ の因子があれば総次数4より $a=4$、残りは全て次数0となる。
次数0の双対から定数束への埋込みは定数なので、切断の値はその固定真部分空間にしか入らない。

$\rho=11$ で上の逆中心 $d$ を用いる。直線 $L$ の $e$ における四平面を $G=W_e$ とする。
$W_p^\perp/G$ は三次元で、$W_d$ との対合で識別される。
$e$ の外の $x\in L$ では $x$ の $T$ への制限は固定された非零方向の倍であり、
$d$ の $L_3$ 鉛筆から $w(d,x)$ は階数3である。
従って四列の $G$ を法とした射影は、$e$ で消える一次式 $t$ と定数階数3行列の積となる。
その核の定数元を $s$ とし、$e$ で非零の補座標を $u$ とする。$s$ に対応する列は
$u e_0+t e_1$ と書ける。全文字で中支持が4なのでこの列には零点がなく、$e_0,e_1\in G$ は独立である。
$e$ での評価が $G$ への同型であることを使い、$s$ の三補元を $e_1,e_2,e_3$ に写るよう選ぶ。
残り三列の $t$ 係数 $h_1,h_2,h_3$ は $G$ を法として独立である。
こうして残りの列は $u e_j+t h_j$ ($j=1,2,3$)。
最初の二列は三行二列のKronecker束、残り二列は各々二行一列の束をなし、
余核は $\mathcal O(2)\oplus2\mathcal O(1)$ である。
\end{proof}
